\chapter{The Signed Archimedean Kernel}
\label{v4-arith:w6-a:ATP-direct-obstruction}

The gamma term in the Weil formula has density
$k(u)=\log\pi-\Re\psi_\Gamma(1/4+iu/2)$.
It is positive near zero and negative at large frequency.
This sign is fixed by Fourier inversion, independently of any
positivity conjecture for the complete arithmetic form.

\section{The two digamma kernels}
\label{v4-arith:w6-a:context}
\label{v4-arith:w6-a:sources}

\begin{definition}[Fourier convention and gamma density]
\label{v4-arith:w6-a:def:g}
For $f\in\mathscr A=C_c^\infty(\mathbb R)$ put
$h_f(u)=\int f(t)e^{iut}dt$. Define
\begin{equation}
 k(u)=\log\pi-\Re\psi_\Gamma(1/4+iu/2),\qquad
 g(u)=\Re\psi_\Gamma(1/4+iu/2)+\log\pi=2\log\pi-k(u).
 \label{v4-arith:w6-a:eq:g-definition}
\end{equation}
The auxiliary form $Q_g$ and the arithmetic gamma form satisfy
\begin{equation}
 \begin{split}
 Q_g(f)&=\frac1{2\pi}\int |h_f(u)|^2g(u)du,\\
 W_\infty(f*f^*)&=\frac1{2\pi}\int |h_f(u)|^2k(u)du
              =2\log\pi\,\|f\|_2^2-Q_g(f).
 \end{split}
 \label{v4-arith:w6-a:eq:W-arch-as-g}
\end{equation}
Both integrals converge because the densities have logarithmic growth.
\end{definition}

\begin{lemma}[digamma series and monotonicity]
\label{v4-arith:w6-a:lem:g-hurwitz}
\label{v4-arith:w6-a:lem:g-monotone}
\label{v4-arith:w6-a:prop:g-boundary}
With $c_n=n+1/4$, one has
\[
 g(u)=-\gamma-\frac\pi2-3\log2+\log\pi
       +\sum_{n=0}^\infty\frac{u^2/4}{c_n(c_n^2+u^2/4)}.
\]
The series and its derivative converge uniformly on compact intervals.
The function $g$ is even and strictly increasing on $(0,\infty)$.
The function $k$ is even and strictly decreasing there, with
\[
 k(0)=\log\pi+\gamma+\pi/2+3\log2>0,\qquad
 k(u)=\log(2\pi/|u|)+O(|u|^{-2}).
\]
\end{lemma}
\begin{proof}
Subtract the digamma series at $1/4$ from its real part at
$1/4+iu/2$. The difference is the displayed positive series.
The value at $1/4$ follows from the reflection and duplication
identities for $\Gamma$. These formulas are
\href{https://dlmf.nist.gov/5.5.E4}{DLMF 5.5.4} and
\href{https://dlmf.nist.gov/5.5.E8}{5.5.8}.
For $|u|\le T$, the summands are $O_T((n+1)^{-3})$.
Differentiation gives
\[
 g'(u)=\sum_{n=0}^\infty
       \frac{c_nu/2}{(c_n^2+u^2/4)^2}>0\quad(u>0).
\]
The differentiated series has the same locally summable bound.
The digamma asymptotic in
\href{https://dlmf.nist.gov/5.11.E2}{DLMF 5.11.2}, in sectors away
from the negative real axis, gives the stated logarithmic asymptotic.
\end{proof}

\begin{proposition}[distinct sign thresholds]
\label{v4-arith:w6-a:thm:t-star}
\label{v4-arith:w6-a:prop:t-star-numerical}
There are unique positive numbers $t_*$ and $u_*$ such that
$g(t_*)=0$ and $k(u_*)=0$. They satisfy $t_*<u_*$.
The gamma density $k$ is positive for $|u|<u_*$ and negative for
$|u|>u_*$. Its threshold is not $t_*$.
\end{proposition}
\begin{proof}
The value $g(0)=-\gamma-\pi/2-3\log2+\log\pi$ is negative:
$\log\pi<\pi/2$ and $\gamma,\log2>0$.
The preceding lemma gives strict monotonicity and divergence to
positive infinity. The intermediate value theorem proves the
existence and uniqueness of $t_*$.
The same argument applies to $k$, which decreases from a positive
value to negative infinity. Since $g(u_*)=2\log\pi>0$, one has
$t_*<u_*$.
\end{proof}

\section{The exact quadratic decomposition}
\label{v4-arith:w6-a:decomposition}

\begin{definition}[the auxiliary core and tail]
\label{v4-arith:w6-a:def:core-tail}
Set
\[
 A[f]=\int_{|u|<t_*}|h_f(u)|^2(-g(u))du,\qquad
 B[f]=\int_{|u|>t_*}|h_f(u)|^2g(u)du.
\]
For the arithmetic kernel itself, set
\[
 C[f]=\int_{|u|<u_*}|h_f(u)|^2k(u)du,\qquad
 D[f]=\int_{|u|>u_*}|h_f(u)|^2(-k(u))du.
\]
These four quantities are nonnegative.
\end{definition}

\begin{proposition}[signed decomposition]
\label{v4-arith:w6-a:lem:decomposition}
\[
 Q_g(f)=\frac{B[f]-A[f]}{2\pi},\qquad
 W_\infty(f*f^*)=\frac{C[f]-D[f]}{2\pi}
 =2\log\pi\,\|f\|_2^2+\frac{A[f]-B[f]}{2\pi}.
\]
In particular, $B[f]\ge A[f]$ characterizes $Q_g(f)\ge0$.
It does not characterize $W_\infty(f*f^*)\ge0$.
\end{proposition}
\begin{proof}
Split the two integrals at their respective sign thresholds and
use Plancherel in \eqref{v4-arith:w6-a:eq:W-arch-as-g}.
\end{proof}

\begin{definition}[frequency dominance conditions]
\label{v4-arith:w6-a:def:HFD-class}
Put $M=-g(0)>0$. The auxiliary high-frequency condition is
\begin{equation}
 2t_*M\sup_{|u|\le t_*}|h_f(u)|^2\le B[f].
 \label{v4-arith:w6-a:eq:HFD-condition}
\end{equation}
It implies $A[f]\le B[f]$, because $-g(u)\le M$ in the core.
The positive gamma cone is $\{f\in\mathscr A:C[f]\ge D[f]\}$.
\end{definition}

\begin{theorem}[both signs on compact smooth squares]
\label{v4-arith:w6-a:thm:HFD-ATP}
\label{v4-arith:w6-a:rem:HFD-nonempty}
For every nonzero $b\in\mathscr A$, the functions
$f_L(t)=L^{-1}b(t/L)$ have $W_\infty(f_L*f_L^*)>0$ for sufficiently
large $L$. The functions $b_v(t)=e^{-ivt}b(t)$ have
$W_\infty(b_v*b_v^*)<0$ for sufficiently large positive $v$.
They satisfy the auxiliary high-frequency condition for large $v$.
For $u_0\in(t_*,u_*)$, the functions
$g_L(t)=L^{-1}b(t/L)e^{-iu_0t}$ have both
$Q_g(g_L)>0$ and $W_\infty(g_L*g_L^*)>0$ for large $L$.
\end{theorem}
\begin{proof}
Since $h_{f_L}(u)=h_b(Lu)$, change of variables gives
\[
 L W_\infty(f_L*f_L^*)
 =\frac1{2\pi}\int |h_b(y)|^2k(y/L)dy
 \longrightarrow k(0)\|b\|_2^2>0.
\]
For $L\ge1$, logarithmic growth bounds the integrand by a fixed
integrable Schwartz weight, so dominated convergence applies.
For modulation, $h_{b_v}(u)=h_b(u-v)$, and
$k(v+y)/\log v\to-1$ for every fixed $y$.
For $v\ge2$, the bound
$|k(v+y)|/\log v\le C(1+\log(2+|y|))$ permits dominated convergence.
Thus
\[
 \frac{W_\infty(b_v*b_v^*)}{\log v}\longrightarrow-\|b\|_2^2.
\]
Equation~\eqref{v4-arith:w6-a:eq:W-arch-as-g} then makes $Q_g(b_v)$
positive for large $v$. More precisely, $B[b_v]$ grows like
$2\pi\|b\|_2^2\log v$, whereas the core supremum decays faster
than every power of $v$. This proves the stronger high-frequency
condition as stated. Finally,
$h_{g_L}(u)=h_b(L(u-u_0))$. The same dominated convergence gives
$LQ_g(g_L)\to g(u_0)\|b\|_2^2>0$ and
$LW_\infty(g_L*g_L^*)\to k(u_0)\|b\|_2^2>0$.
\end{proof}

\section{The complete Weil inequality}
\label{v4-arith:w6-a:bounded-interval}
\label{v4-arith:w6-a:bounded-interval:ruled-out}

Put
\begin{equation}
 R[f]=P(f*f^*)-W_{\mathrm{fin}}(f*f^*).
 \label{v4-arith:w6-a:eq:R-definition}
\end{equation}
The polar form has both signs. The finite-prime form also has both
signs by Proposition~\ref{v4-arith:sc:two-signs}.

\begin{theorem}[the arithmetic inequality with all terms retained]
\label{v4-arith:w6-a:thm:bounded-reduction}
\label{v4-arith:w6-a:thm:irreducible-ATP}
For each $f\in\mathscr A$, the inequality $W(f*f^*)\ge0$ is
exactly either of the following equivalent inequalities:
\begin{equation}
 C[f]\le D[f]+2\pi R[f],\qquad
 A[f]+4\pi\log\pi\,\|f\|_2^2\le B[f]+2\pi R[f].
 \label{v4-arith:w6-a:eq:bounded-reduction-star}
\end{equation}
Requiring these inequalities for all $f\in\mathscr A$ is equivalent
to RH. The separate positive gamma cone does not remove the terms
in $R[f]$.
\end{theorem}
\begin{proof}
Insert $W=R-W_\infty$ and the signed decomposition.
The equivalence to RH is Weil's criterion on $\mathscr A$.
\end{proof}

\begin{proposition}[limits of frequency restrictions]
\label{v4-arith:w6-a:prop:rule-out-projection}
\label{v4-arith:w6-a:prop:rule-out-majorant}
\label{v4-arith:w6-a:prop:rule-out-pw-shrinkage}
No nonzero $f\in\mathscr A$ has $h_f$ supported in a bounded real
frequency interval. Neither compact support of $f$ nor nonnegativity
of the auxiliary form implies a fixed sign of the gamma form.
\end{proposition}
\begin{proof}
The Fourier transform $h_f$ is entire. Vanishing on a real interval
forces it to vanish identically, and Fourier inversion gives $f=0$.
The dilation and modulation examples prove the sign assertions.
\end{proof}

\begin{remark}[comparison with the complete arithmetic form]
\label{v4-arith:w6-a:rem:live-constructions}
The exact signed representation
\eqref{v4-arith:sc:full-comparison} places the gamma term, both polar
moments, and the prime translations on a common test algebra.
A positive realization of that complete form requires the full
inequality \eqref{v4-arith:w6-a:eq:bounded-reduction-star}.
\end{remark}
